
\documentclass{article}
\usepackage{colortbl}
\usepackage{makecell}
\usepackage{multirow}
\usepackage{supertabular}

\begin{document}

\newcounter{utterance}

\centering \large Interaction Transcript for game `codenames', experiment `frequency\_high', episode 8 with reflexion{-}llama{-}3.1{-}8b{-}instruct.
\vspace{24pt}

{ \footnotesize  \setcounter{utterance}{1}
\setlength{\tabcolsep}{0pt}
\begin{supertabular}{c@{$\;$}|p{.15\linewidth}@{}p{.15\linewidth}p{.15\linewidth}p{.15\linewidth}p{.15\linewidth}p{.15\linewidth}}
    \# & $\;$A & \multicolumn{4}{c}{Game Master} & $\;\:$B\\
    \hline

    \theutterance \stepcounter{utterance}  
    & & & \multicolumn{2}{p{0.3\linewidth}}{
        \cellcolor[rgb]{0.9,0.9,0.9}{
            \makecell[{{p{\linewidth}}}]{
                \texttt{\tiny{[GM$|$GM]}}
                {'hidden': {'team': ['future', 'message'], 'innocent': ['lunch', 'car', 'control', 'blood', 'dinner'], 'opponent': [], 'assassin': ['state']}, 'revealed': {'team': {'team': [], 'innocent': [], 'opponent': [], 'assassin': []}, 'opponent': {'team': [], 'innocent': [], 'opponent': [], 'assassin': []}}}
            }
        }
    }
    & & \\ \\

    \theutterance \stepcounter{utterance}  
    & & \multicolumn{4}{p{0.6\linewidth}}{
        \cellcolor[rgb]{0.9,0.9,0.9}{
            \makecell[{{p{\linewidth}}}]{
                \texttt{\tiny{[P1$\langle$GM]}}
                \texttt{You have attempted this same task before and did not succeed. These are the plans you wrote after those attempts.} \\
\texttt{Trial \#1: New plan: In my last attempt, I provided a clue that related to multiple words, but I didn't consider the fact that one of the related words was an opponent's word, which caused my teammate's turn to end prematurely. Next time, I will make sure to carefully review the lists of words and avoid including any opponent's words in my clue. I will also consider providing a clue that relates to fewer words, but is more specific and easier to guess. For example, I could provide a clue that relates to only one or two words, such as "CLUE: future TARGETS: future, charge".} \\
\\ 
\texttt{Trial \#2: New plan: In my last attempt, I provided a clue that related to multiple words, but I didn't consider the fact that one of the related words was an opponent's word, which caused my teammate's turn to end prematurely. However, this time I also made another mistake: I provided a clue that was too broad and didn't exclude the distractor words. Next time, I will make sure to carefully review the lists of words and avoid including any opponent's words in my clue. I will also consider providing a clue that relates to fewer words, but is more specific and easier to guess. Additionally, I will make sure to exclude distractor words from my clue. For example, I could provide a clue that relates to only one or two words, such as "CLUE: future TARGETS: future, charge", and also exclude distractor words like "lunch" or "car".} \\
\\ 
\texttt{Trial \#3: New plan: In my last attempts, I provided clues that related to multiple words, but I didn't consider the fact that one of the related words was an opponent's word, which caused my teammate's turn to end prematurely. I also didn't exclude distractor words from my clue. Next time, I will make sure to carefully review the lists of words and avoid including any opponent's words in my clue. I will also consider providing a clue that relates to fewer words, but is more specific and easier to guess. Additionally, I will make sure to exclude distractor words from my clue. To achieve this, I will carefully examine the lists of words and identify the most relevant and specific connections between the team words and the remaining words. I will then provide a clue that is narrow in scope, yet still provides a clear and useful hint for my teammate. For example, I could provide a clue that relates to only one word, such as "CLUE: future TARGETS: future", and also exclude distractor words like "lunch" or "car".} \\
\\ 
\\ 
\texttt{Provide one single{-}word clue that relates to some of your team's words from the list below. You can choose to how many and to which words your clue relates to.} \\
\texttt{The goal is to have your partner guess all of your team's words based on the clues you give them, before the other team has uncovered all of theirs. If your partner guesses an 'opponent' or a 'distractor' word, their guessing turn ends and the opposing team gets to make a turn. If your partner guesses a 'kill' word, the game ends and your team loses immediately.} \\
\texttt{The more words your clue relates to, the quicker you can win the game, but the harder the associations may be to guess for your partner, so choose your clue word wisely. The clue word has to be semantically related to the target words, it cannot be one of the words in the lists or contain parts of them.} \\
\\ 
\texttt{Always give your single{-}word clue and your comma{-}separated list of related target words in the following format and make your answers as short as possible, never include any other text than is required in this form:} \\
\\ 
\texttt{CLUE: <WORD>} \\
\texttt{TARGETS: <TARGETS>} \\
\\ 
\texttt{Your team words are: future, message.} \\
\texttt{Your opponent's team words are: .} \\
\texttt{Distractor words are: lunch, car, control, blood, dinner.} \\
\texttt{The kill words are: state.} \\
            }
        }
    }
    & \\ \\

    \theutterance \stepcounter{utterance}  
    & \multicolumn{4}{p{0.6\linewidth}}{
        \cellcolor[rgb]{0.9,0.9,0.9}{
            \makecell[{{p{\linewidth}}}]{
                \texttt{\tiny{[P1$\rangle$GM]}}
                \texttt{CLUE: plan} \\
\texttt{TARGETS: future, message} \\
            }
        }
    }
    & & \\ \\

    \theutterance \stepcounter{utterance}  
    & & & \multicolumn{2}{p{0.3\linewidth}}{
        \cellcolor[rgb]{0.9,0.9,0.9}{
            \makecell[{{p{\linewidth}}}]{
                \texttt{\tiny{[GM$|$GM]}}
                \texttt{plan} \\
            }
        }
    }
    & & \\ \\

    \theutterance \stepcounter{utterance}  
    & & & \multicolumn{2}{p{0.3\linewidth}}{
        \cellcolor[rgb]{0.9,0.9,0.9}{
            \makecell[{{p{\linewidth}}}]{
                \texttt{\tiny{[GM$|$GM]}}
                ['future', 'message']
            }
        }
    }
    & & \\ \\

    \theutterance \stepcounter{utterance}  
    & & & \multicolumn{2}{p{0.3\linewidth}}{
        \cellcolor[rgb]{0.9,0.9,0.9}{
            \makecell[{{p{\linewidth}}}]{
                \texttt{\tiny{[GM$|$GM]}}
                {'word': 'future', 'assignment': 'team'}
            }
        }
    }
    & & \\ \\

    \theutterance \stepcounter{utterance}  
    & & & \multicolumn{2}{p{0.3\linewidth}}{
        \cellcolor[rgb]{0.9,0.9,0.9}{
            \makecell[{{p{\linewidth}}}]{
                \texttt{\tiny{[GM$|$GM]}}
                {'word': 'message', 'assignment': 'team'}
            }
        }
    }
    & & \\ \\

    \theutterance \stepcounter{utterance}  
    & & & \multicolumn{2}{p{0.3\linewidth}}{
        \cellcolor[rgb]{0.9,0.9,0.9}{
            \makecell[{{p{\linewidth}}}]{
                \texttt{\tiny{[GM$|$GM]}}
                \texttt{opponent has won} \\
            }
        }
    }
    & & \\ \\

    \theutterance \stepcounter{utterance}  
    & \multicolumn{4}{p{0.6\linewidth}}{
        \cellcolor[rgb]{0.95,0.95,0.95}{
            \makecell[{{p{\linewidth}}}]{
                \texttt{\tiny{[P2$\rightarrow$P1]}}
                \texttt{Post{-}episode Reflexion (not sent to GM):} \\
\texttt{New plan: In my last attempts, I provided clues that related to multiple words, but I didn't consider the fact that one of the related words was an opponent's word, which caused my teammate's turn to end prematurely. I also didn't exclude distractor words from my clue. Next time, I will make sure to carefully review the lists of words and avoid including any opponent's words in my clue. I will also consider providing a clue that relates to fewer words, but is more specific and easier to guess. Additionally, I will make sure to exclude distractor words from my clue. To achieve this, I will carefully examine the lists of words and identify the most relevant and specific connections between the team words and the remaining words. I will then provide a clue that is narrow in scope, yet still provides a clear and useful hint for my teammate. Furthermore, I will also consider the kill words and make sure that my clue does not inadvertently lead to the game ending prematurely. For example, I could provide a clue that relates to only one word, such as "CLUE: future TARGETS: future", and also exclude distractor words like "lunch" or "car".} \\
            }
        }
    }
    & & \\ \\

\end{supertabular}
}

\end{document}
